% GENERATED FILE. Edit canonical lessons, then run npm run generate:books.
% canonical-source-hash: fnv1a64:b94a01d66de98f22
% canonical-lessons: ZH-C348-taolun, ZH-C348-shangliang, ZH-C348-zhenglun, ZH-C348-duifang, ZH-C348-yizhi

\chapter{Talking It Over}
\label{ch:zh-talking-it-over}
\hlchaptermodality{348}

\begin{chapteropening}
\textbf{By the end of this chapter:} \emph{I can say we are discussing something, talk it over with the other side and say we all agree.}

Complete the last lesson of chapter 348: I can say we are discussing something, talk it over with the other side and say we all agree.
\end{chapteropening}

\section[tǎolùn]{tǎolùn — to discuss}

\label{lesson:ZH-C348-taolun}



\begin{quote}
Before the new one: say the Mandarin for the significance of something, then the Mandarin for to know from experience.
\end{quote}

\subsection*{You'll want to know: tǎolùn}
\emph{tǎolùn} — \textquotedblleft{}to discuss\textquotedblright{}.

\begin{quote}
\emph{wǒmen míngtiān kāihuì tǎolùn zhège wèntí} — We'll discuss this at tomorrow's meeting.
\end{quote}

Its characters come later; for now it is a word for the ear.

\begin{cousinweb}
\emph{tǎolùn} is built from \emph{tǎo} and \emph{lùn}, both inherited from Middle Chinese, \emph{thawX} and \emph{lwon} (Guangyun, in Baxter's transcription), which glosses the first 'to govern' and 'to punish', and the second 'to speak' and 'to discuss'; it also lists the second in the departing tone, the tone of lùn today.
\end{cousinweb}

\begin{grammarlens}[title={What you've built}]
One more word for giving reasons.
\end{grammarlens}

\subsection*{Your turn}
\begin{itemize}
\raggedright
  \item \emph{Say it:} \emph{tǎolùn}
  \item \emph{Say it:} \emph{tǎolùn}, once more
  \item \emph{Say it:} \emph{wǒmen míngtiān kāihuì tǎolùn zhège wèntí} — We'll discuss this at tomorrow's meeting.
  \item \emph{Recall:} say the Mandarin for the significance of something, then the Mandarin for to know from experience, then say \emph{tǎolùn} again
\end{itemize}

\begin{tcolorbox}[breakable,colback=teal!4,colframe=teal!35!black,title={Before you move on}]
What does tǎolùn mean? (\textquotedblleft{}To discuss\textquotedblright{}.) Say it once more.
\end{tcolorbox}
\section[shāngliang]{shāngliang — to talk it over together; to consult}

\label{lesson:ZH-C348-shangliang}



\begin{quote}
Before the new one: say the Mandarin for to know from experience, then the Mandarin for to discuss.
\end{quote}

\subsection*{You'll want to know: shāngliang}
\emph{shāngliang} — \textquotedblleft{}to talk it over together; to consult\textquotedblright{}.

\begin{quote}
\emph{wǒmen yīnggāi gēn lǎoshī shāngliang} — We should talk it over with the teacher.
\end{quote}

Its characters come later; for now it is a word for the ear.

\begin{cousinweb}
\emph{shāngliang} is built from \emph{shāng} and \emph{liáng}, both inherited from Middle Chinese, \emph{syang} and \emph{ljang} (Guangyun, in Baxter's transcription), which glosses the first, among other senses, 'to deliberate', and the second 'to measure'. The second is said light here.
\end{cousinweb}

\begin{grammarlens}[title={What you've built}]
One more word for giving reasons.
\end{grammarlens}

\subsection*{Your turn}
\begin{itemize}
\raggedright
  \item \emph{Say it:} \emph{shāngliang}
  \item \emph{Say it:} \emph{shāngliang}, once more
  \item \emph{Say it:} \emph{wǒmen yīnggāi gēn lǎoshī shāngliang} — We should talk it over with the teacher.
  \item \emph{Recall:} say the Mandarin for to know from experience, then the Mandarin for to discuss, then say \emph{shāngliang} again
\end{itemize}

\begin{tcolorbox}[breakable,colback=teal!4,colframe=teal!35!black,title={Before you move on}]
What does shāngliang mean? (\textquotedblleft{}To talk it over together; to consult\textquotedblright{}.) Say it once more.
\end{tcolorbox}
\section[zhēnglùn]{zhēnglùn — to argue over a point; to debate}

\label{lesson:ZH-C348-zhenglun}



\begin{quote}
Before the new one: say the Mandarin for to discuss, then the Mandarin for to talk it over together.
\end{quote}

\subsection*{You'll want to know: zhēnglùn}
\emph{zhēnglùn} — \textquotedblleft{}to argue over a point; to debate\textquotedblright{}.

\begin{quote}
\emph{tāmen zài zhēnglùn zhège wèntí} — They're arguing over this question.
\end{quote}

Its characters come later; for now it is a word for the ear.

\begin{cousinweb}
\emph{zhēnglùn} is built from \emph{zhēng} and \emph{lùn}, both inherited from Middle Chinese, \emph{tsreang} and \emph{lwon} (Guangyun, in Baxter's transcription), which glosses the first 'to compete' and the second 'to discuss'.
\end{cousinweb}

\begin{grammarlens}[title={What you've built}]
One more word for giving reasons.
\end{grammarlens}

\subsection*{Your turn}
\begin{itemize}
\raggedright
  \item \emph{Say it:} \emph{zhēnglùn}
  \item \emph{Say it:} \emph{zhēnglùn}, once more
  \item \emph{Say it:} \emph{tāmen zài zhēnglùn zhège wèntí} — They're arguing over this question.
  \item \emph{Recall:} say the Mandarin for to discuss, then the Mandarin for to talk it over together, then say \emph{zhēnglùn} again
\end{itemize}

\begin{tcolorbox}[breakable,colback=teal!4,colframe=teal!35!black,title={Before you move on}]
What does zhēnglùn mean? (\textquotedblleft{}To argue over a point; to debate\textquotedblright{}.) Say it once more.
\end{tcolorbox}
\section[duìfāng]{duìfāng — the other side; the other party}

\label{lesson:ZH-C348-duifang}



\begin{quote}
Before the new one: say the Mandarin for to talk it over together, then the Mandarin for to argue over a point.
\end{quote}

\subsection*{You'll want to know: duìfāng}
\emph{duìfāng} — \textquotedblleft{}the other side; the other party\textquotedblright{}.

\begin{quote}
\emph{nǐ yīnggāi tīng duìfāng shuō wán} — You should hear the other side out.
\end{quote}

Its characters come later; for now it is a word for the ear.

\begin{cousinweb}
\emph{duìfāng} is built from \zh{对} 'facing' and \emph{fāng} 'a side', both inherited from Middle Chinese, \emph{twojH} and \emph{pjang} (Guangyun, in Baxter's transcription), which glosses the first 'to answer' and 'to match', and the second 'the four quarters'.
\end{cousinweb}

\begin{grammarlens}[title={What you've built}]
One more word for giving reasons.
\end{grammarlens}

\subsection*{Your turn}
\begin{itemize}
\raggedright
  \item \emph{Say it:} \emph{duìfāng}
  \item \emph{Say it:} \emph{duìfāng}, once more
  \item \emph{Say it:} \emph{nǐ yīnggāi tīng duìfāng shuō wán} — You should hear the other side out.
  \item \emph{Recall:} say the Mandarin for to talk it over together, then the Mandarin for to argue over a point, then say \emph{duìfāng} again
\end{itemize}

\begin{tcolorbox}[breakable,colback=teal!4,colframe=teal!35!black,title={Before you move on}]
What does duìfāng mean? (\textquotedblleft{}The other side; the other party\textquotedblright{}.) Say it once more.
\end{tcolorbox}
\section[yízhì]{yízhì — all of one mind; in agreement}

\label{lesson:ZH-C348-yizhi}



\begin{quote}
Before the new one: say the Mandarin for to argue over a point, then the Mandarin for the other party.
\end{quote}

\subsection*{You'll want to know: yízhì}
\emph{yízhì} — \textquotedblleft{}all of one mind; in agreement\textquotedblright{}.

\begin{quote}
\emph{dàjiā de yìjiàn hěn yízhì} — Everyone was of one mind.
\end{quote}

Its characters come later; for now it is a word for the ear.

The tone of yī changes here: before a falling tone it is said yí, and the spelling writes the change.

\begin{cousinweb}
\emph{yízhì} is built from \zh{一} 'one' and \emph{zhì} 'to arrive', both inherited from Middle Chinese, \emph{'jit} and \emph{trijH} (Guangyun, in Baxter's transcription), which glosses the first 'the start of the numbers' and the second 'to arrive'. How the pair came to mean agreement is not settled.
\end{cousinweb}

\begin{grammarlens}[title={What you've built}]
One more word for giving reasons.
\end{grammarlens}

\subsection*{Your turn}
\begin{itemize}
\raggedright
  \item \emph{Say it:} \emph{yízhì}
  \item \emph{Say it:} \emph{yízhì}, once more
  \item \emph{Say it:} \emph{dàjiā de yìjiàn hěn yízhì} — Everyone was of one mind.
  \item \emph{Recall:} say the Mandarin for to argue over a point, then the Mandarin for the other party, then say \emph{yízhì} again
\end{itemize}

\begin{tcolorbox}[breakable,colback=teal!4,colframe=teal!35!black,title={Before you move on}]
What does yízhì mean? (\textquotedblleft{}All of one mind; in agreement\textquotedblright{}.) Say it once more.
\end{tcolorbox}
